% =====================================================================
% Le Collé — Équations différentielles : exercices complémentaires du TD
% =====================================================================

\begin{exo}{edn-1}{1}{premier ordre, second membre simple}
Résoudre sur $\R$ l'équation $y'+2ty=t$. Déterminer la solution qui vérifie $y(0)=0$.
\end{exo}
\begin{corrige}{edn-1}
Homogène : une primitive de $2t$ est $t^2$, d'où $y=Ke^{-t^2}$. La constante $\frac12$ est solution évidente ($0+2t\cdot\frac12=t$).
Solutions : $y=\frac12+Ke^{-t^2}$ ; $y(0)=0$ donne $K=-\frac12$, soit $y=\frac12\big(1-e^{-t^2}\big)$.
\gain{toujours tester une constante quand le second membre est proportionnel au coefficient de $y$.}
\end{corrige}

\begin{exo}{edn-2}{2}{normaliser avant de résoudre}
Résoudre sur $\R_+^*$ l'équation $2ty'+y=1+t$.
\end{exo}
\begin{corrige}{edn-2}
Forme normalisée : $y'+\frac1{2t}y=\frac{1+t}{2t}$. Une primitive de $\frac1{2t}$ est $\frac12\ln t$ : homogène $y=Ke^{-\frac12\ln t}=\frac K{\sqrt t}$.
On cherche une solution affine $y=a+bt$ : $2bt+a+bt=1+t$, d'où $a=1$, $b=\frac13$. Solutions : $y=1+\frac t3+\frac K{\sqrt t}$.
\end{corrige}

\begin{exo}{edn-3}{2}{racine double et second membre polynôme-exponentielle}
Résoudre sur $\R$ l'équation $y''-4y'+4y=(t+1)e^{2t}$.
\end{exo}
\begin{corrige}{edn-3}
$r^2-4r+4=(r-2)^2$ : racine double $2$, $S_H=\{(A+Bt)e^{2t}\}$. Le second membre est $P(t)e^{2t}$ avec $2$ racine double :
on cherche $y=Q(t)e^{2t}$ et, comme $m^2+am+b=0$ et $2m+a=0$, il reste $Q''=t+1$ : $Q=\frac{t^3}6+\frac{t^2}2$ (on peut prendre
les constantes nulles). Solutions : $y=\left(\frac{t^3}6+\frac{t^2}2+A+Bt\right)e^{2t}$.
\retenir{racine double : $Q''=P$, il suffit d'intégrer deux fois.}
\end{corrige}

\begin{exo}{edn-4}{2}{problème de Cauchy, racines complexes}
Déterminer la solution réelle de $y''+2y'+5y=0$ telle que $y(0)=1$ et $y'(0)=1$, et décrire son comportement en $+\infty$.
\end{exo}
\begin{corrige}{edn-4}
$r^2+2r+5=0$ : $r=-1\pm2i$. Donc $y=e^{-t}(A\cos2t+B\sin2t)$. $y(0)=A=1$ ; $y'(0)=-A+2B=1$, $B=1$.
$y=e^{-t}(\cos2t+\sin2t)=\sqrt2\,e^{-t}\cos\left(2t-\frac\pi4\right)$ : oscillations de pulsation $2$, amorties, qui tendent vers $0$.
\end{corrige}

\begin{exo}{edn-5}{2}{superposition}
Résoudre sur $\R$ l'équation $y''-3y'+2y=\sh(2t)$.
\end{exo}
\begin{corrige}{edn-5}
Racines $1$ et $2$ : $S_H=\{Ae^t+Be^{2t}\}$. On écrit $\sh(2t)=\frac12e^{2t}-\frac12e^{-2t}$ et l'on superpose.
Pour $\frac12e^{2t}$ ($m=2$ racine simple) : $Q''+(2m+a)Q'+(m^2+am+b)Q=Q''+Q'=\frac12$, $Q=\frac t2$.
Pour $-\frac12e^{-2t}$ ($m=-2$ non racine) : $(4+6+2)Q=-\frac12$, $Q=-\frac1{24}$.
Solutions : $y=\frac t2e^{2t}-\frac1{24}e^{-2t}+Ae^t+Be^{2t}$.
\piege{$\sh(2t)$ n'est pas « de la forme » d'une seule exponentielle : on le décompose avant de chercher.}
\end{corrige}

\begin{exo}{edn-6}{3}{recollement : une famille à deux paramètres}
On considère $(E)$ : $ty'-2y=t^3$.
\begin{enumerate}
\item Résoudre $(E)$ sur $\R_+^*$ et sur $\R_-^*$.
\item Déterminer toutes les solutions de $(E)$ sur $\R$. Le théorème de Cauchy est-il contredit ?
\end{enumerate}
\end{exo}
\begin{corrige}{edn-6}
\q{1.} Sur $\R_+^*$ ou $\R_-^*$ : $y'-\frac2ty=t^2$. Homogène : $e^{2\ln|t|}=t^2$, donc $Kt^2$ ; $y=t^3$ convient
($3t^3-2t^3=t^3$). Solutions : $y=t^3+K_\pm t^2$ sur chaque demi-droite.

\q{2.} Soit $y$ solution sur $\R$. Les deux expressions tendent vers $0$ en $0$ : $y(0)=0$. Les taux d'accroissement
$\frac{y(t)}t=t^2+K_\pm t$ tendent tous deux vers $0$ : $y$ est dérivable en $0$, $y'(0)=0$, \emph{quels que soient} $K_+$ et $K_-$.
L'équation en $0$ s'écrit $0-0=0$ : elle est vérifiée. Les solutions sur $\R$ sont donc les fonctions égales à $t^3+K_+t^2$ pour
$t\geq0$ et à $t^3+K_-t^2$ pour $t\leq0$, avec $K_+$, $K_-$ arbitraires.
Ce n'est pas une contradiction : le coefficient de $y'$ s'annule en $0$, et le théorème de Cauchy ne s'applique que sur un
intervalle où l'équation se normalise.
\retenir{au point de raccord, il peut y avoir aucune, une seule ou une infinité de solutions.}
\end{corrige}

\begin{exo}{edn-7}{3}{développement limité d'une solution}
Soit $f$ la solution sur $]-\infty,1[$ de $2(t-1)y'+y=\sin(2t)+t^2$ telle que $f(0)=0$. Sans résoudre l'équation, déterminer le
développement limité de $f$ à l'ordre $4$ en $0$.
\end{exo}
\begin{corrige}{edn-7}
\choix{on n'a pas besoin de $f$ : l'équation fournit une relation entre les coefficients du développement.}
Sur $]-\infty,1[$, $f'=\frac{\sin2t+t^2-f}{2(t-1)}$ ; par récurrence, $f$ est de classe $\mathcal C^\infty$, donc admet un
développement $f(t)=\sum_{k=0}^5a_kt^k+o(t^5)$, et $f'(t)=\sum_{k=1}^5ka_kt^{k-1}+o(t^4)$. Le coefficient de $t^k$ dans
$2tf'-2f'+f$ est $(2k+1)a_k-2(k+1)a_{k+1}$, et $\sin2t+t^2=2t+t^2-\frac43t^3+o(t^4)$. Avec $a_0=f(0)=0$ :
\[a_0-2a_1=0,\quad3a_1-4a_2=2,\quad5a_2-6a_3=1,\quad7a_3-8a_4=-\tfrac43,\]
d'où $a_1=0$, $a_2=-\frac12$, $a_3=-\frac7{12}$, $a_4=-\frac{11}{32}$ :
\[f(t)=-\frac{t^2}2-\frac{7t^3}{12}-\frac{11t^4}{32}+o(t^4).\]
\end{corrige}

\begin{exo}{edn-8}{3}{résonance}
Soit $\omega>0$. Résoudre $y''+\omega^2y=\cos t$, en distinguant $\omega\neq1$ et $\omega=1$. Dans quel cas toutes les solutions
sont-elles bornées ?
\end{exo}
\begin{corrige}{edn-8}
$S_H=\{A\cos\omega t+B\sin\omega t\}$. On résout $z''+\omega^2z=e^{it}$.
\begin{itemize}
\item $\omega\neq1$ : $i$ n'est pas racine de $r^2+\omega^2$, $z=\frac{e^{it}}{\omega^2-1}$, d'où $y_p=\frac{\cos t}{\omega^2-1}$.
Toutes les solutions sont bornées.
\item $\omega=1$ : $i$ est racine simple, $z=Q(t)e^{it}$ avec $Q''+2iQ'=1$, $Q=-\frac i2t$ ; $y_p=\operatorname{Re}z=\frac t2\sin t$.
Aucune solution n'est bornée (le terme $\frac t2\sin t$ ne l'est pas).
\end{itemize}
\retenir{excitation à la pulsation propre : résonance, amplitude non bornée.}
\end{corrige}

\begin{exo}{edn-9}{3}{variation de la constante au second ordre}
Résoudre sur $\left]-\frac\pi2,\frac\pi2\right[$ l'équation $y''+y=\dfrac1{\cos t}$.
\end{exo}
\begin{corrige}{edn-9}
\choix{le second membre n'est pas de la forme $P(t)e^{mt}$ : variation de la constante, avec une solution de $(H)$ qui ne s'annule pas sur l'intervalle, $\varphi=\cos$.}
Avec $y=\lambda\cos t$ : $y''+y=\lambda''\cos t-2\lambda'\sin t$, d'où, pour $\mu=\lambda'$ : $\mu'-2\tan(t)\,\mu=\frac1{\cos^2t}$.
Homogène : $\mu=\frac K{\cos^2t}$ ; variation : $\mu=\frac{C(t)}{\cos^2t}$ avec $C'=1$. Donc $\lambda'=\frac{t+K}{\cos^2t}$ et, par parties,
$\lambda=(t+K)\tan t+\ln(\cos t)+D$. Solutions :
\[y=t\sin t+\cos t\,\ln(\cos t)+K\sin t+D\cos t.\]
\end{corrige}

\begin{exo}{edn-10}{4}{équation d'Euler}
Résoudre sur $\R_+^*$ l'équation $t^2y''-3ty'+4y=t^2$ en posant $t=e^x$.
\end{exo}
\begin{corrige}{edn-10}
$x\mapsto e^x$ est une bijection $\mathcal C^2$ de $\R$ sur $\R_+^*$, de réciproque $\ln$. Avec $z(x)=y(e^x)$ :
$z'=e^xy'(e^x)$, $z''=e^{2x}y''(e^x)+e^xy'(e^x)$, donc $t^2y''=z''-z'$ et $ty'=z'$. L'équation devient
$z''-4z'+4z=e^{2x}$ : racine double $2$, $Q''=1$, $z=\left(\frac{x^2}2+A+Bx\right)e^{2x}$. En revenant à $t$ :
\[y=t^2\left(\frac{(\ln t)^2}2+A+B\ln t\right).\]
\retenir{équation d'Euler $t^2y''+aty'+by$ : $t=e^x$ donne $t^2y''=z''-z'$, $ty'=z'$.}
\end{corrige}

\begin{exo}{edn-11}{4}{série entière et fonctions hyperboliques}
On considère $(E)$ : $4ty''+2y'-y=0$.
\begin{enumerate}
\item Déterminer les solutions de $(E)$ développables en série entière, et leur rayon de convergence.
\item Exprimer la somme à l'aide de fonctions usuelles sur $\R_+$, puis sur $\R_-$.
\end{enumerate}
\end{exo}
\begin{corrige}{edn-11}
\q{1.} Avec $y=\sum a_nt^n$ : le coefficient de $t^n$ est $4(n+1)na_{n+1}+2(n+1)a_{n+1}-a_n=2(n+1)(2n+1)a_{n+1}-a_n$. Donc
$a_{n+1}=\frac{a_n}{(2n+2)(2n+1)}$ et $a_n=\frac{a_0}{(2n)!}$. Le rapport $\left|\frac{a_{n+1}}{a_n}\right|$ tend vers $0$ : rayon infini.
Les solutions développables en série entière sont $y=a_0\sum\frac{t^n}{(2n)!}$.

\q{2.} Pour $t\geq0$, $\sum\frac{t^n}{(2n)!}=\sum\frac{(\sqrt t)^{2n}}{(2n)!}=\ch\sqrt t$ ; pour $t\leq0$, avec $t=-u^2$,
$\sum\frac{(-1)^nu^{2n}}{(2n)!}=\cos\sqrt{-t}$.
\gain{reconnaître $(2n)!$ au dénominateur : ce sont les termes pairs de $\exp$, $\ch$ ou $\cos$.}
\end{corrige}

\begin{exo}{edn-12}{4}{changement de fonction et recollement}
On considère $(E)$ : $t^2y''+4ty'+(2-t^2)y=1$.
\begin{enumerate}
\item On pose $z(t)=t^2y(t)$. Montrer que $y$ est solution de $(E)$ sur un intervalle si et seulement si $z''-z=1$.
\item Résoudre $(E)$ sur $\R_+^*$ et sur $\R_-^*$.
\item Montrer que $(E)$ admet une unique solution sur $\R$.
\end{enumerate}
\end{exo}
\begin{corrige}{edn-12}
\q{1.} $z'=2ty+t^2y'$ et $z''=2y+4ty'+t^2y''$, donc $t^2y''+4ty'+(2-t^2)y=z''-z$.

\q{2.} $z''-z=1$ : $z=-1+A\ch t+B\sh t$. Sur chaque demi-droite, $y=\frac{-1+A_\pm\ch t+B_\pm\sh t}{t^2}$.

\q{3.} Une solution sur $\R$ est continue en $0$. Quand $t\to0^+$, $-1+A_+\ch t+B_+\sh t=(A_+-1)+B_+t+O(t^2)$ : la limite de $y$ n'est
finie que si $A_+=1$ et $B_+=0$ ; de même à gauche. Il reste $y(t)=\frac{\ch t-1}{t^2}$ pour $t\neq0$, et $y(0)=\frac12$ (limite).
Cette fonction est la somme de la série entière $\sum\frac{t^{2n}}{(2n+2)!}$, de rayon infini : elle est de classe
$\mathcal C^\infty$ sur $\R$, et l'équation en $0$ s'écrit $2\cdot\frac12=1$. Il y a donc une unique solution sur $\R$.
\end{corrige}

\begin{exo}{edn-13}{4}{une solution polynomiale, puis l'autre}
On considère $(E)$ : $(1-t^2)y''-2ty'+2y=0$ sur $]-1,1[$.
\begin{enumerate}
\item Déterminer les solutions polynomiales de $(E)$.
\item En posant $y=t\,\lambda(t)$ sur $]0,1[$, montrer que $\lambda'$ vérifie une équation du premier ordre, et la résoudre.
\item En déduire que les solutions de $(E)$ sur $]-1,1[$ sont les fonctions $t\mapsto at+b\left(\dfrac t2\ln\dfrac{1+t}{1-t}-1\right)$.
\end{enumerate}
\end{exo}
\begin{corrige}{edn-13}
\q{1.} Coefficient de $t^n$ : $\big(-n(n-1)-2n+2\big)a_n=-(n-1)(n+2)a_n$ : $n=1$. Avec $y=at+b$ : $-2at+2at+2b=0$, $b=0$. Solutions polynomiales : $y=at$.

\q{2.} $y=t\lambda$ : $y'=\lambda+t\lambda'$, $y''=2\lambda'+t\lambda''$. En reportant, les termes en $\lambda$ disparaissent :
$(1-t^2)t\lambda''+\big(2(1-t^2)-2t^2\big)\lambda'=0$. Pour $\mu=\lambda'$ : $\mu'=-\frac{2-4t^2}{t(1-t^2)}\mu$. La décomposition
$\frac{2-4t^2}{t(1-t^2)}=\frac2t-\frac1{1-t}+\frac1{1+t}$ donne $\mu=\frac K{t^2(1-t^2)}=K\left(\frac1{t^2}+\frac1{2(1-t)}+\frac1{2(1+t)}\right)$.

\q{3.} En intégrant : $\lambda=K\left(-\frac1t+\frac12\ln\frac{1+t}{1-t}\right)+a$, et $y=t\lambda=at+K\left(\frac t2\ln\frac{1+t}{1-t}-1\right)$.
Cette fonction est définie et deux fois dérivable sur $]-1,1[$ tout entier et l'on vérifie qu'elle est solution (ou on invoque
l'unicité de Cauchy de part et d'autre de $0$). Les deux solutions $t$ et $\frac t2\ln\frac{1+t}{1-t}-1$ ne sont pas
proportionnelles : elles engendrent le plan $S_H$.
\retenir{une solution connue $\varphi$ (même si elle s'annule en un point) : $y=\lambda\varphi$ ramène au premier ordre sur les intervalles où $\varphi\neq0$.}
\end{corrige}

\begin{exo}{edn-14}{5}{problème : trois méthodes pour une même équation}
On considère $(E)$ : $ty''+2y'+ty=0$.
\begin{enumerate}
\item Déterminer les solutions de $(E)$ développables en série entière au voisinage de $0$, et reconnaître leur somme.
\item On pose $z(t)=ty(t)$. Montrer que $y$ est solution de $(E)$ sur un intervalle si et seulement si $z''+z=0$.
\item En déduire les solutions de $(E)$ sur $\R_+^*$ et sur $\R_-^*$.
\item Déterminer les solutions de $(E)$ sur $\R$. Forment-elles un plan vectoriel ? Pourquoi le théorème de Cauchy ne
s'applique-t-il pas ?
\end{enumerate}
\end{exo}
\begin{corrige}{edn-14}
\q{1.} Avec $y=\sum a_nt^n$ : $ty''=\sum(n+1)na_{n+1}t^n$, $2y'=\sum2(n+1)a_{n+1}t^n$, $ty=\sum a_{n-1}t^n$ ($n\geq1$). Le coefficient
constant donne $2a_1=0$ ; pour $n\geq1$, $(n+1)(n+2)a_{n+1}+a_{n-1}=0$. Les termes impairs sont nuls, et
$a_{2p}=\frac{(-1)^pa_0}{(2p+1)!}$ (récurrence). Rayon infini, et $y=a_0\sum\frac{(-1)^pt^{2p}}{(2p+1)!}=a_0\frac{\sin t}t$ ($t\neq0$),
prolongée par $a_0$ en $0$.

\q{2.} $z'=y+ty'$, $z''=2y'+ty''$, donc $z''+z=ty''+2y'+ty$.

\q{3.} $z=A\cos t+B\sin t$, donc $y=\frac{A\cos t+B\sin t}t$ sur chaque demi-droite (constantes $A_\pm$, $B_\pm$).

\q{4.} Pour qu'une solution ait une limite finie en $0$, il faut $A_+=A_-=0$ (sinon $\frac{A\cos t}t$ diverge). Il reste
$y=B_\pm\frac{\sin t}t$, de limite $B_\pm$ : la continuité impose $B_+=B_-=B$. La fonction $B\frac{\sin t}t$ (prolongée par $B$) est la
somme d'une série entière de rayon infini, donc $\mathcal C^\infty$, et elle est solution sur $\R$ (question~1). Les solutions sur
$\R$ forment une \emph{droite} vectorielle, et non un plan : le coefficient $t$ de $y''$ s'annule en $0$, l'équation ne se
normalise pas sur $\R$ et le théorème de Cauchy ne s'y applique pas (par exemple, aucune solution ne vérifie $y(0)=0$, $y'(0)=1$).
\retenir{trois méthodes (série entière, changement de fonction, recollement) qui se recoupent et se contrôlent mutuellement.}
\end{corrige}

\begin{exo}{edn-15}{5}{solutions périodiques d'un oscillateur forcé}
Soit $f$ une fonction continue sur $\R$ et l'équation $(E)$ : $y''+y=f(t)$.
\begin{enumerate}
\item Montrer que $y_0(t)=\displaystyle\int_0^t\sin(t-u)f(u)\,du$ est la solution de $(E)$ telle que $y_0(0)=y_0'(0)=0$.
\emph{(Développer $\sin(t-u)$ pour dériver.)}
\item On suppose $f$ $2\pi$-périodique, et l'on pose $C=\int_0^{2\pi}f(u)\cos u\,du$, $S=\int_0^{2\pi}f(u)\sin u\,du$. Montrer que
$y_0(t+2\pi)-y_0(t)=C\sin t-S\cos t$.
\item En déduire que $(E)$ admet une solution $2\pi$-périodique si et seulement si $C=S=0$, et qu'alors toutes les solutions
le sont.
\item Application : $f(t)=\cos t$, puis $f(t)=\cos(2t)$.
\end{enumerate}
\end{exo}
\begin{corrige}{edn-15}
\q{1.} $y_0(t)=\sin t\int_0^tf(u)\cos u\,du-\cos t\int_0^tf(u)\sin u\,du$. En dérivant (produit et intégrale de borne variable),
les termes $\sin t\,f(t)\cos t-\cos t\,f(t)\sin t$ s'annulent : $y_0'(t)=\cos t\int_0^tf\cos+\sin t\int_0^tf\sin$. En dérivant
encore : $y_0''=-y_0+f(t)(\cos^2t+\sin^2t)=-y_0+f$. Enfin $y_0(0)=y_0'(0)=0$.

\q{2.} Comme $f(u)\cos u$ est $2\pi$-périodique, $\int_0^{t+2\pi}f\cos=\int_0^tf\cos+\int_t^{t+2\pi}f\cos=\int_0^tf\cos+C$, et de même avec
$\sin$ ; en reportant dans l'expression de $y_0$, on obtient $y_0(t+2\pi)=y_0(t)+C\sin t-S\cos t$.

\q{3.} Les solutions sont $y=y_0+A\cos t+B\sin t$, et $y(t+2\pi)-y(t)=C\sin t-S\cos t$ pour toutes. Cette fonction est nulle si et
seulement si $C=S=0$ (prendre $t=0$ et $t=\frac\pi2$). Donc ou bien toutes les solutions sont $2\pi$-périodiques, ou bien aucune.

\q{4.} $f=\cos$ : $C=\pi\neq0$, aucune solution périodique (résonance : on retrouve $\frac t2\sin t$). $f(t)=\cos2t$ : $C=S=0$, toutes
les solutions sont $2\pi$-périodiques (par exemple $y_p=-\frac13\cos2t$).
\end{corrige}

\begin{exo}{edn-16}{1}{deux seconds membres}
Résoudre sur $\R$ l'équation $y''+y=t+e^t$.
\end{exo}
\begin{corrige}{edn-16}
$S_H=\{A\cos t+B\sin t\}$ (racines $\pm i$). Par superposition : $y=t$ convient pour le second membre $t$, et $y=ce^t$ donne
$2c=1$ pour $e^t$. Solutions : $y=t+\frac12e^t+A\cos t+B\sin t$.
\end{corrige}

\begin{exo}{edn-17}{2}{variation de la constante au premier ordre}
Résoudre sur $\R$ l'équation $(1+t^2)y'+ty=1$.
\end{exo}
\begin{corrige}{edn-17}
Forme normalisée : $y'+\frac t{1+t^2}y=\frac1{1+t^2}$ ; une primitive de $\frac t{1+t^2}$ est $\frac12\ln(1+t^2)$, d'où
$S_H=\left\{\frac K{\sqrt{1+t^2}}\right\}$. Aucune solution évidente : on pose $y=\frac{K(t)}{\sqrt{1+t^2}}$, et il reste
$\frac{K'(t)}{\sqrt{1+t^2}}=\frac1{1+t^2}$, soit $K'=\frac1{\sqrt{1+t^2}}$, $K=\argsh t+C$. Solutions :
\[y=\frac{\argsh t+C}{\sqrt{1+t^2}}=\frac{\ln\big(t+\sqrt{1+t^2}\big)+C}{\sqrt{1+t^2}}.\]
\retenir{en variation de la constante, les termes en $K$ se simplifient toujours : il ne reste que $K'\varphi=$ second membre.}
\end{corrige}

\begin{exo}{edn-18}{3}{une équation fonctionnelle}
Déterminer les fonctions $f$ dérivables sur $\R$ telles que, pour tout réel $x$, $f'(x)=f(-x)$.
\end{exo}
\begin{corrige}{edn-18}
\choix{on se ramène à une équation différentielle en dérivant la relation.}
Soit $f$ une solution. Comme $x\mapsto f(-x)$ est dérivable, $f'$ l'est : $f$ est deux fois dérivable et
$f''(x)=-f'(-x)=-f(x)$. Donc $f=A\cos+B\sin$. Réciproquement, $f'(x)=-A\sin x+B\cos x$ et $f(-x)=A\cos x-B\sin x$ :
l'égalité pour tout $x$ équivaut à $B=A$. Les solutions sont les fonctions $x\mapsto A(\cos x+\sin x)$.
\piege{l'équation $f''=-f$ n'est qu'une condition \emph{nécessaire} : la réciproque élimine une constante.}
\end{corrige}
